% 1.引言.tex —— 一、引言
\section{引言}

\subsection{问题背景}
干燥是决定中药材成品品质的关键工序之一。热风烘干通过调控烘房温湿环境完成药材干燥，其工艺参数选取不当会导致干燥效率低、能耗高、成品品质不稳定。传统试验优化模式成本高、周期长，亟需借助数理分析与数值仿真得到干燥规律。

预热平衡阶段是干燥的第一阶段，药材被烘房加热、内部温度逐步趋近环境、表层水分开始向表面迁移并汽化。该阶段内部温度与含水率的分布直接决定了后续恒温干燥阶段的初始状态与工艺时点的选取，因此需要求解一个含变系数、带对流边界的瞬态热质传递问题。

\subsection{问题重述}
\textbf{问题 1}：某中药材形状大致呈圆柱形，长 $25$ cm、半径 $2$ cm。烘干开始时药材温度为 $28$ $^\circ$C、水分浓度（干基含水率）为 $2.55$ kg/kg，烘房温度与水分浓度随时间变化见附件 1。请建立预热平衡阶段药材温度和水分浓度变化规律的数学模型（相关参数见附录 2），在论文中按表 1、表 2 格式给出 $100,300,600,900,1200,1500,1800$ s、到药材中心距离 $0,0.5,1,1.5,2$ cm 处的温度与水分浓度，并将 $1800$ s 内每隔 $1$ s、每隔 $0.1$ cm 的完整结果保存到 \texttt{result1.xlsx}。所有结果保留四位小数。
